% ============================================================
% 卷首语
% ============================================================
\colchapter{卷首语}

\section*{我们都是对比下长大的孩子}
\auth{文/木子懿}

放假回家的时候，跟着家里人和街坊邻居去吃夜宵。大人们之间的聊天主题总是工作和小孩的成绩。邻居阿姨笑眯眯地看着我说：“要是妹妹也像你一样会读书就好了，我现在天天在家告诉妹妹让她向你学习，不要天天想着出去玩，多在家里学习，以后考个大学。”

我乖巧地对阿姨笑了笑，没有说话。她家的孩子我认识，一个学习挺认真的女孩，今年还考上了市重点高中。我看着妈妈和阿姨开心地讨论着育儿经验，摆在桌上的夜宵被微凉的夜风吹着，夜宵店的背景音乐混杂着人们的讨论声，把我的心绪扯得很乱。

好像每一个孩子都逃不过被对比的命运，成长的每一步有一点偏差就需要用更加完美的人来矫正，大人会摁着你一步步靠近别人。

小学时，有一次我好不容易借到了同学的漫画全集，因为有太多同学排队要看，我必须在三天内看完，于是晚上偷偷摸摸地打着台灯躲在被子里翻阅。不知道是翻书的声响还是被戳中笑点发出了声音，最终被爸爸发现了。他怒不可遏地质问我：“你表姐在被子里开灯看的都是语文书，你看看你看的是什么？”的确，那时我的成绩非常一般。

表姐，我对她再了解不过，因为从小爸妈就拿她来教育我，她是所有大人眼里的好学生，年年都有三好学生的奖状。可这一切从我初中毕业后发生了巨大的改变。那一年，表姐高考落榜了，而我却考上了省重点高中。据长辈说，表姐是因为早恋导致的成绩下滑，于是，她不再是优秀的代表，反而成了警示弟弟妹妹们的“反面教材”。

我还有一个双胞胎妹妹，从小到大，不管是样貌还是成绩，无一避免，全都要被拿来对比。妹妹从小成绩就比我好，省重点高中高手如云，有一次妹妹考了班级第四，而我考了班级倒数第四。于是我经常要接受一些其他人的“定论”：我比妹妹蠢，比妹妹懒，没有妹妹努力。还好，我亲爱的爸爸妈妈从来不这样想，他们没有区别对待，更没有放弃我。妹妹也没有看低我，而是一直帮助我。

每个人的成长都离不开一个“别人家的孩子”。经常听大人们说：“你怎么不和××一样？”只是这样简简单单的一句话，就可能会让不自信的种子在我们心里扎根。如果能把“什么都比不上人家”改成“不用什么都跟人家比”，可能我们在成长的道路上就会多一份自信，多一份张扬吧！我一直觉得，每个人都是独一无二的。不管怎样，做好自己，我们就是长大了。
